\documentclass[11pt]{amsart}
\usepackage[margin=1in]{geometry}
\usepackage{amsmath,amssymb,amsthm}
\usepackage{booktabs}
\usepackage[utf8]{inputenc}
\usepackage{newunicodechar}
% Unicode in the verbatim Lean excerpt (pdflatex); xelatex/lualatex with a Unicode mono font also works.
\newunicodechar{ℕ}{\ensuremath{\mathbb{N}}}
\newunicodechar{ℚ}{\ensuremath{\mathbb{Q}}}
\newunicodechar{ℝ}{\ensuremath{\mathbb{R}}}
\newunicodechar{→}{\ensuremath{\to}}
\newunicodechar{⊕}{\ensuremath{\oplus}}
\newunicodechar{∑}{\ensuremath{\Sigma}}
\newunicodechar{≤}{\ensuremath{\le}}
\newunicodechar{∧}{\ensuremath{\wedge}}
\newunicodechar{∃}{\ensuremath{\exists}}
\newunicodechar{∀}{\ensuremath{\forall}}
\usepackage{tabularx}
\usepackage{url}
\usepackage[hidelinks]{hyperref}

\newtheorem{theorem}{Theorem}
\newtheorem{lemma}[theorem]{Lemma}
\newtheorem{corollary}[theorem]{Corollary}
\theoremstyle{remark}
\newtheorem{remark}[theorem]{Remark}

\newcommand{\Q}{\mathbb{Q}}
\newcommand{\Z}{\mathbb{Z}}
\newcommand{\C}{\mathbb{C}}
\newcommand{\cP}{\mathcal{P}}
\newcommand{\cL}{\mathcal{L}}
\newcommand{\ip}[1]{\langle #1\rangle}

\title[The $N^{\sqrt2+o(1)}$ construction for Hurwitz's problem]{The $N^{\sqrt2+o(1)}$ construction for Hurwitz's problem,\\ written out in terms of bilinear forms}
\author{Lindsay Debnath}
\date{Version 1, 10 November 2026}

\begin{document}
\maketitle

\begin{abstract}
For every $N$ there are integer bilinear forms $f_1,g_1,\dots,f_K,g_K$ in $x_1,\dots,x_N,y_1,\dots,y_N$ with
$h!\,(\sum_i x_i^2)(\sum_j y_j^2)=\sum_{k\le K} f_kg_k$ and $K=N^{\sqrt2+o(1)}$, where $h=O(\log^2N)$ is the least
integer with at least $N$ partitions. This is a consequence of
Theorem~B of the paper \emph{A $2^{o(n)}$ time algorithm for 3-SAT} \cite{ETH}, first pointed out in
\cite{Sie}. This note restates the construction directly as an identity of bilinear forms, gives the coefficients
explicitly, writes out the proofs of the two lemmas it rests on with the compressed steps filled in, and says
exactly what is machine-checked. Nothing here is new mathematics beyond \cite{ETH}; the point is to make the
construction checkable by a reader who works with polynomials rather than with the function language of \cite{ETH}.
\end{abstract}

\section{The result and its status}

Write $p(E)$ for the number of partitions of $E$ ($p(0)=1$, $p(E)=0$ for $E<0$), $\cP_h$ for the set of
partitions of $h$, and $B_F(N)$ for the least $s$ with $(\sum_{i\le N}x_i^2)(\sum_{j\le N}y_j^2)=\sum_{k\le s}f_kf'_k$,
$f_k,f'_k$ bilinear over $F$ (the quantity $B_F(\mathrm{SOS}_N)$ of \cite{HWY}, and the bilinear-product variant (9) of \cite{Hru}).

\begin{theorem}\label{thm:main}
Let $h\ge1$ and $K_h=\sum_{j=0}^{2h}(4h-2j+1)\,p(j)$. There are bilinear forms $f_k,g_k$ ($k\le K_h$; some
products $f_kg_k$ may vanish) in
variables $x_\xi,y_\eta$ indexed by $\xi,\eta\in\cP_h$, with integer coefficients given explicitly in
Section~\ref{sec:forms}, such that
\[
  h!\Big(\sum_{\xi\in\cP_h}x_\xi^2\Big)\Big(\sum_{\eta\in\cP_h}y_\eta^2\Big)=\sum_{k=1}^{K_h}f_kg_k .
\]
\end{theorem}

Setting all but $N\le p(h)$ of the $x$'s and $y$'s to zero gives the identity for $N$ variables. Taking $h$ least
with $p(h)\ge N$, the Hardy--Ramanujan asymptotics give $K_h=N^{\sqrt2+o(1)}$ (Section~\ref{sec:count}, for $N=2^\lambda$;
for other $N$ use $2^{\lambda-1}<N\le2^\lambda$). Hence:

\begin{corollary}\label{cor}
$B_\Q(N)\le N^{\sqrt2+o(1)}$, and the same over any field of characteristic $0$ or larger than $h=O(\log^2N)$.
By \cite[Remark~1.5]{HWY}, $\sigma_\C(N)\le 3B_\C(N)\le N^{\sqrt2+o(1)}$, where $\sigma_\C(N)$ is the least number of
squares in Hurwitz's problem over $\C$.
\end{corollary}

The previous best upper bound is $N^{c+o(1)}$ with $c=1+2H(\alpha/2)\approx1.6147$, $H(\alpha)=\tfrac12$
(\cite[Theorem~1 and its proof in \S4]{Hru}, stated there as $O(N^{c'})$ with $c'<1.62$; for the bilinear-product form (9)
see \cite[\S5.1, Remark~18]{Hru}, and \cite[Theorem~12]{HruCCC} for a direct proof over every field). It
also holds in positive characteristic, where the present identity says nothing ($h!=0$ once $p\le h$). The open problem
in \cite[\S7]{Hru} asks whether $\sigma_\C(N)\le N^{1.5+o(1)}$ is attainable by parity representations; the bound
here is below that target, by a different route.

\begin{remark}[a consequence we have not checked beyond substitution]
The proof of \cite[Theorem~25]{Hru} uses Theorem~1 there only as a black box, $\sigma(n)=O(n^{1+\delta})$,
and obtains $\mathrm{size}^{(nc)}(\mathrm{ID}_n)\le n^{(3+\omega(3)+\delta)/4+o(1)}$. Substituting
$\delta=\sqrt2-1$ and $\omega(3)<4.1997$ gives exponent below $1.904$ (instead of $1.954$), over fields of
characteristic $0$.
\end{remark}

\paragraph{Status.} Theorem~B of \cite{ETH} is formalized in Lean~4 with Mathlib in \cite{FM} (commit
\texttt{16a5ce41}): the statement file \texttt{eth/proof/statements/AppHalfPiecesStatement.lean} restates Appendix~B
in the paper's words (definitions and statements only, about 240 lines, importing Mathlib and two short statement
files), and \texttt{GeneralPieces.thmHalfPiecesPure} proves it. We rebuilt this from source
(\texttt{lake exe cache get; lake build GeneralPieces.Main}) and checked
\texttt{\#print axioms}: only \texttt{propext}, \texttt{Classical.choice}, \texttt{Quot.sound}. We then stated
Corollary~\ref{cor} (the bound over $\Q$, for $N$ a power of two) in terms of polynomials, in a file that imports Mathlib only
(excerpt, inside \texttt{namespace HurwitzSqrt2} with \texttt{open MvPolynomial Finset}; docstrings omitted):
{\footnotesize
\begin{verbatim}
abbrev Idx (lam : ℕ) : Type := Fin lam → Bool
abbrev Var (lam : ℕ) : Type := Idx lam ⊕ Idx lam
noncomputable def bil {lam : ℕ} (F : Idx lam → Idx lam → ℚ) : MvPolynomial (Var lam) ℚ :=
  ∑ i, ∑ j, C (F i j) * X (Sum.inl i) * X (Sum.inr j)
noncomputable def sos (lam : ℕ) : MvPolynomial (Var lam) ℚ :=
  (∑ i, X (Sum.inl i) ^ 2) * (∑ j, X (Sum.inr j) ^ 2)
def Statement : Prop :=
  ∃ c : ℕ, ∀ lam : ℕ, 1 ≤ lam →
    ∃ (K : ℕ) (f g : Fin K → Idx lam → Idx lam → ℚ),
      (K : ℝ) ≤ (2 * lam : ℝ) ^ c * 2 ^ (Real.sqrt 2 * lam) ∧
      sos lam = ∑ k, bil (f k) * bil (g k)
\end{verbatim}
}
and proved \texttt{Statement} from \texttt{thmHalfPiecesPure} (file \texttt{HurwitzProof.lean}, 292 lines). That file
restates the \texttt{HurwitzSqrt2} block rather than importing it; we diffed the two blocks and they are identical, and
\texttt{\#print axioms hurwitzSqrt2} gives the same three axioms. The Lean statement has rational coefficients and does
not contain the factor $h!$, integrality, or the exact count $K_h$ (its $K$ is
$\sum_s2^{\lceil\log_2\min(R^{\rm row}_s,R^{\rm col}_s)\rceil}\le2K_h$); those are checked numerically for $h\le5$
(Section~\ref{sec:map}) and follow from (54), (55).
So the claim rests on Lean's kernel plus a reading of the definitions above and of Definition~7 as formalized
(\texttt{Pieces.Family.Is} in \texttt{PiecesStatement.lean}, which matches the printed definition field by field,
except the running-time clause of 7(c), which the bound does not need).
The Lean development was written and reviewed by AI models; no person is known to have read its proofs.
Section~\ref{sec:map} maps each step below to its Lean declaration.

\section{The construction}\label{sec:construction}

\subsection{The Fock space}
Let $V=\Q[t_1,t_2,\dots]$, with $t_k$ of weight $k$; the monomials $t^\mu=\prod_it_{\mu_i}$, $\mu$ a partition,
form a basis, and $V_E$ (weight $E$) has dimension $p(E)$. For $k\ge1$ let
\[\nabla_{-k}=t_k\cdot,\qquad \nabla_k=k\,\partial/\partial t_k,\qquad\text{so}\quad[\nabla_k,\nabla_{-l}]=k\,\delta_{kl}.\]
These are the Heisenberg relations; $V$ is the Fock representation. $\nabla_k$ maps $V_E\to V_{E-k}$ for all $k\ne0$.
For $\sigma,\sigma'$ partitions and an operator $T$, write $[t^{\sigma'}]\,T\,t^\sigma$ for the matrix entry.

\subsection{The operators $a[r]$}
For a monomial $a=t^\mu$ with parts $\mu_1,\dots,\mu_n$ (listed with multiplicity) and $r\in\Z$, put
\begin{equation}\tag{47}
a[r]=\sum_{\substack{k_1,\dots,k_n\in\Z\setminus\{0\}\\ k_1+\dots+k_n=-r}}\;\prod_i\binom{-k_i-1}{\mu_i-1}\;
\prod_{i:k_i<0}\nabla_{k_i}\prod_{i:k_i>0}\nabla_{k_i},
\end{equation}
with all differentiations ($k_i>0$) applied first, and extend linearly in $a$. Here
$\binom{k}{j}=k(k-1)\cdots(k-j+1)/j!$ for all integers $k$. On $V_E$ only finitely many terms are nonzero (the
positive $k_i$ add up to at most $E$), and $a[r]$ maps $V_E$ into $V_{E+r}$, with integer matrix entries if $a$ has
integer coefficients. For example $t_1[r]=\nabla_{-r}$ ($r\neq0$). With $Y(a,\omega)=\sum_ra[r]\,\omega^{r-|\mu|}$, these are the
coefficients of the fields of the Heisenberg vertex algebra \cite[\S2.3, (4)]{Fre}, as \cite[App.~H]{ETH} notes; we do
not use this except in Remark~\ref{rem:locality}.

\subsection{The operators $Q_s$}
Let $h\ge1$, $b\in V_h$, and $a=\sum_{h'=0}^h\bar a_{h'}$ with $\bar a_{h'}\in V_{h'}$. Put
$a^{\langle i\rangle}=\sum_{h'}\binom{h'+h}{i}\bar a_{h'}$ and, for $s\in\Z$,
\begin{equation}\tag{48}
Q_s(a,b)=\sum_{i=0}^{2h}(-1)^i\,a^{\langle i\rangle}[s-h+i]\;b[h-s-i],\qquad
Q'_s(a,b)=\sum_{i=0}^{2h}(-1)^i\,b[h-s-i]\;a^{\langle i\rangle}[s-h+i].
\end{equation}
Both preserve weight. The weighting $a^{\langle i\rangle}$ makes $Q_s$ linear in $a$ while each homogeneous component
$\bar a_{h'}$ is paired with its own binomial row $\binom{h'+h}{i}$; this is what lets the proof below handle
inhomogeneous $a$ one component at a time.

\begin{lemma}[Lemma 9 of \cite{ETH}]\label{lem:Q}
On $V_E$, for every $E\ge0$:
(a) $Q_s(a,b)=Q'_s(a,b)$;
(b) $Q_s(a,b)=0$ for $|s|>E+h$;
(c) $\sum_{|s|\le E+h}Q_s(a,b)=(a[-h]b)\cdot\mathrm{Id}$, where $a[-h]b\in V_0=\Q$.
\end{lemma}

\subsection{Dual families}
For $n\ge1$ let $y_n=\sum_{j=1}^n\binom nj t_j$ (no relation to the variables $y_\eta$ of Theorem~\ref{thm:main}), and for $\mu\in\cP_h$
\[a^\mu=\prod_i(-1)^{\mu_i-1}y_{\mu_i},\qquad b^\mu=\frac{h!}{z_\mu}\,t^\mu,\qquad z_\mu=\prod_kk^{m_k(\mu)}m_k(\mu)!,\]
where $m_k(\mu)$ is the number of parts equal to $k$; $h!/z_\mu\in\Z$ is the number of permutations of cycle type $\mu$.
$a^\mu$ has integer coefficients and only monomials of weight $\le h$.

\begin{lemma}[Lemma 10 of \cite{ETH}]\label{lem:dual}
$a^\mu[-h]\,b^\nu=h!\,\delta_{\mu\nu}$ for all $\mu,\nu\in\cP_h$.
\end{lemma}

\subsection{The identity}
For $\xi,\eta,\upsilon,\nu\in\cP_h$ and $|s|\le2h$ let
\begin{equation}\tag{53}
A_s(\xi,\eta,\upsilon,\nu)=[t^\upsilon]\;Q_s(a^\xi,b^\nu)\;t^\eta .
\end{equation}
By Lemma~\ref{lem:Q}(b),(c) with $E=h$ and Lemma~\ref{lem:dual},
\begin{equation}\label{eq:cross}
\sum_{|s|\le2h}A_s(\xi,\eta,\upsilon,\nu)=h!\,\delta_{\xi\nu}\,\delta_{\eta\upsilon}.
\end{equation}
Equivalently, the degree-four polynomial $\sum_s\sum A_s(\xi,\eta,\upsilon,\nu)\,x_\xi y_\eta x_\nu y_\upsilon$ equals
$h!(\sum x_\xi^2)(\sum y_\eta^2)$, already at the level of coefficient tensors (so also for the non-commutative
$\mathrm{ID}$ polynomial of \cite{HWY}, with $(\xi,\eta,\nu,\upsilon)$ in the positions $X,Y,X,Y$). What remains is to
write each tensor $A_s$ as a short sum of products of two matrices, splitting the indices as $(\xi,\upsilon)\,|\,(\nu,\eta)$
or as $(\xi,\eta)\,|\,(\nu,\upsilon)$; commutatively, each such product is a product of two bilinear forms. (These two
splittings are forms (b) and (a) of \cite[Lemma~3.5]{HWY}, which matters for the non-commutative width of
$h!\,\mathrm{ID}_N$; Theorem~\ref{thm:main} itself does not need that lemma.)

\section{The bilinear forms}\label{sec:forms}

Define integer matrices (rows and columns indexed by partitions) for each $\xi,\nu\in\cP_h$:
\[
\alpha^{(i)}_r[\xi]_{\sigma'\sigma}=(-1)^i\,[t^{\sigma'}]\;(a^\xi)^{\langle i\rangle}[r]\;t^\sigma,\qquad
\beta_r[\nu]_{\sigma'\sigma}=[t^{\sigma'}]\;b^\nu[r]\;t^\sigma .
\]
Inserting $\mathrm{Id}=\sum_\sigma t^\sigma[t^\sigma]$ on the middle weight space in (48) gives two factorizations.

\emph{Column-wise} (directly from $Q_s$; display (55) of \cite{ETH}): $b[h-s-i]\,t^\eta\in V_{2h-s-i}$, so
\[
A_s(\xi,\eta,\upsilon,\nu)=\sum_{i=0}^{2h}\sum_{\sigma\in\cP_{2h-s-i}}\alpha^{(i)}_{s-h+i}[\xi]_{\upsilon\sigma}\;\beta_{h-s-i}[\nu]_{\sigma\eta},
\]
and each term $(i,\sigma)$ gives the pair
\[
f_{s,i,\sigma}=\sum_{\xi,\upsilon}\alpha^{(i)}_{s-h+i}[\xi]_{\upsilon\sigma}\,x_\xi y_\upsilon,\qquad
g_{s,i,\sigma}=\sum_{\nu,\eta}\beta_{h-s-i}[\nu]_{\sigma\eta}\,x_\nu y_\eta .
\]

\emph{Row-wise} (from $Q'_s$, so it uses Lemma~\ref{lem:Q}(a); display (54)): $a^{\langle i\rangle}[s-h+i]\,t^\eta\in V_{s+i}$, so
\[
A_s(\xi,\eta,\upsilon,\nu)=\sum_{i=0}^{2h}\sum_{\sigma\in\cP_{s+i}}\alpha^{(i)}_{s-h+i}[\xi]_{\sigma\eta}\;\beta_{h-s-i}[\nu]_{\upsilon\sigma},
\]
\[
f_{s,i,\sigma}=\sum_{\xi,\eta}\alpha^{(i)}_{s-h+i}[\xi]_{\sigma\eta}\,x_\xi y_\eta,\qquad
g_{s,i,\sigma}=\sum_{\nu,\upsilon}\beta_{h-s-i}[\nu]_{\upsilon\sigma}\,x_\nu y_\upsilon .
\]

In both cases $f_{s,i,\sigma}g_{s,i,\sigma}$, summed over the terms, has coefficient $A_s(\xi,\eta,\upsilon,\nu)$ on
$x_\xi y_\eta x_\nu y_\upsilon$, so by \eqref{eq:cross} the sum over $s$ is $h!(\sum x^2)(\sum y^2)$.

\paragraph{Which way for which $s$.} The numbers of terms are $R^{\rm row}_s=\sum_{i=0}^{2h}p(s+i)$ and
$R^{\rm col}_s=\sum_{i=0}^{2h}p(2h-s-i)$. Substituting $i\mapsto 2h-i$ shows $R^{\rm col}_s=R^{\rm row}_{-s}$, and
$R^{\rm row}_s$ is nondecreasing in $s$. So the column-wise form is at least as short for $s\ge0$ and the row-wise
form for $s\le0$, and
\[
\min(R^{\rm row}_s,R^{\rm col}_s)=\sum_{j=0}^{2h-|s|}p(j),\qquad
K_h=\sum_{|s|\le2h}\ \sum_{j=0}^{2h-|s|}p(j)=\sum_{j=0}^{2h}(4h-2j+1)\,p(j).
\]
This proves Theorem~\ref{thm:main}, given Lemmas~\ref{lem:Q} and~\ref{lem:dual}. Note that Lemma~\ref{lem:Q}(a)
is used only for $s<0$. The intuition: both factors of a term pass through one weight space, of weight $2h-s-i$
or $s+i$; choosing the side that keeps it at most $2h-|s|$ keeps its dimension at most $p(2h)\approx p(h)^{\sqrt2}$.

\section{Proofs of the lemmas}\label{sec:proofs}

\begin{proof}[Proof of Lemma~\ref{lem:dual}]
For $b\in V_h$ and a monomial $a$, in $a[-h]b$ only terms of (47) with all $k_i>0$ survive: if some $k_i<0$, the
positive ones add up to more than $h$ and the differentiations kill $b$. Expand $\prod_iy_{\mu_i}=\sum_{j}\prod_i\binom{\mu_i}{j_i}t_{j_i}$
($1\le j_i\le\mu_i$) and apply (47) to each monomial, listing its parts as $j_1,\dots,j_n$:
\[
\Big(\prod_iy_{\mu_i}\Big)[-h]\,b=\sum_{\substack{k_i\ge1\\ \sum k_i=h}}\prod_i\vartheta_{\mu_i}(k_i)\prod_i\nabla_{k_i}b,\qquad
\vartheta_n(k)=\sum_{j=1}^n\binom{n}{n-j}\binom{-k-1}{j-1}=\binom{n-k-1}{n-1}
\]
by Vandermonde's identity (valid for all integer upper indices). For $1\le k<n$, $0\le n-k-1<n-1$ and $\vartheta_n(k)=0$;
$\vartheta_n(n)=\binom{-1}{n-1}=(-1)^{n-1}$. So a nonzero term has $k_i\ge\mu_i$ for all $i$; as both sides add up to $h$,
$k_i=\mu_i$. The signs cancel those in $a^\mu$, so $a^\mu[-h]b=\prod_i\nabla_{\mu_i}b$. For $b=t^\nu$,
$\prod_i\mu_i\partial_{t_{\mu_i}}t^\nu$ is zero unless $\mu\subseteq\nu$ as multisets, i.e.\ $\mu=\nu$, and then
equals $z_\mu$. Hence $a^\mu[-h]b^\nu=\frac{h!}{z_\nu}z_\mu\delta_{\mu\nu}=h!\,\delta_{\mu\nu}$.
\end{proof}

\begin{proof}[Proof of Lemma~\ref{lem:Q}]
All three statements are linear in $a$ and $b$ (since $a^{\langle i\rangle}$ is linear in $a$), so it suffices to
treat monomials. Let $a=t^\mu$, $|\mu|=h'\le h$, $b=t^\nu$, $|\nu|=h$, and put $H=h'+h$; then
$a^{\langle i\rangle}=\binom Hi a$. The terms of (48) with $i>H$ vanish, so (48) is a sum over $0\le i\le H$, which is
what the expansion of $(\omega_1-\omega_2)^H$ below produces.

\emph{Rings.} Let $\cL=\Q[\omega_1^{\pm1},\omega_2^{\pm1}]$; $\cL^+$ the formal series $\sum_{j\ge j_0}q_j(\omega_1)\omega_2^j$ with
$q_j\in\Q[\omega_1^{\pm1}]$, and $\cL^-$ the same with $\omega_1,\omega_2$ exchanged. In $\cL^\pm$, $\omega_1-\omega_2$ has
the inverse $\varepsilon_+=\sum_{m\ge0}\omega_1^{-m-1}\omega_2^m$, resp.\ $\varepsilon_-=-\sum_{m\ge0}\omega_2^{-m-1}\omega_1^m$.
From $(\omega_1-\omega_2)\varepsilon_\pm=1$: $\partial_{\omega_1}\varepsilon_\pm=-\varepsilon_\pm^2$, $\partial_{\omega_2}\varepsilon_\pm=\varepsilon_\pm^2$.
Hence $\partial_{\omega_2}^{n-1}\varepsilon^2=n!\,\varepsilon^{n+1}$ and
\begin{equation}\tag{49}
\frac{1}{(j-1)!(n-1)!}\,\partial_{\omega_1}^{j-1}\partial_{\omega_2}^{n-1}\varepsilon_\pm^2=\Gamma_{j,n}\,\varepsilon_\pm^{j+n},\qquad
\Gamma_{j,n}=(-1)^{j-1}\frac{(j+n-1)!}{(j-1)!\,(n-1)!},
\end{equation}
the same number for both signs.

\emph{Generating functions.} For $n\ge1$ let
\[J^+_n(\omega)=\sum_{l\ge1}\binom{l-1}{n-1}\nabla_{-l}\,\omega^{l-n},\qquad
J^-_n(\omega)=\sum_{k\ge1}\binom{-k-1}{n-1}\nabla_k\,\omega^{-k-n}.\]
Reading (47) with $k_i=-l_i<0$ for $i$ in a set $S$ of parts and $k_i>0$ otherwise,
\[Y(t^\mu,\omega):=\sum_r t^\mu[r]\,\omega^{r-h'}=\sum_{S}\prod_{i\in S}J^+_{\mu_i}(\omega)\prod_{i\notin S}J^-_{\mu_i}(\omega),\]
creation operators to the left. Take commuting indeterminates $z_n,z'_n$ ($n\ge1$),
$A^\pm=\sum_nz_nJ^\pm_n(\omega_1)$, $B^\pm=\sum_nz'_nJ^\pm_n(\omega_2)$. The $J^+$'s commute with each other, as do the
$J^-$'s, so expanding the exponentials and collecting $z^\mu=\prod z_{\mu_i}$,
\[e^{A^+}e^{A^-}=\sum_\mu \frac{z^\mu}{\mu!}\,Y(t^\mu,\omega_1),\qquad \mu!=\prod_km_k(\mu)!,\]
because the coefficient of $z^\mu$ collects the terms with $\alpha_n+\beta_n=m_n(\mu)$ factors $J^+_n$, $J^-_n$;
and $1/(\alpha_n!\beta_n!)=\binom{m_n}{\alpha_n}/m_n!$, where $\binom{m_n}{\alpha_n}$ counts the ways to choose which
$\alpha_n$ of the $m_n$ parts equal to $n$ lie in $S$. Likewise for $B^\pm$ with $\omega_2$.

\emph{Commuting the middle factors.} $B^+$ is multiplication by $\sum_kB_kt_k$, $B_k=\sum_nz'_n\binom{k-1}{n-1}\omega_2^{k-n}$,
and $A^-=\sum_kA_k\nabla_k$, $A_k=\sum_jz_j\binom{-k-1}{j-1}\omega_1^{-k-j}$. By Taylor's formula, on the coefficient of each power of $\omega_2$,
$e^{A^-}$ acts as the ring homomorphism $t_k\mapsto t_k+kA_k$, so
$e^{A^-}e^{B^+}=e^{\kappa}e^{B^+}e^{A^-}$ with
\[
\kappa=\sum_kkA_kB_k=\sum_{j,n}z_jz'_n\sum_{k\ge1}k\binom{-k-1}{j-1}\binom{k-1}{n-1}\omega_1^{-k-j}\omega_2^{k-n}
=\sum_{j,n}z_jz'_n\,\Gamma_{j,n}\,\varepsilon_+^{j+n}\quad\text{in }\cL^+,
\]
since the inner sum is the left side of (49) applied to $\varepsilon_+^2=\partial_{\omega_2}\varepsilon_+=\sum_kk\,\omega_1^{-k-1}\omega_2^{k-1}$.
In the same way $e^{B^-}e^{A^+}=e^{\kappa'}e^{A^+}e^{B^-}$ in $\cL^-$, $\kappa'=\sum z_jz'_n\Gamma_{j,n}\varepsilon_-^{j+n}$
(now using $\varepsilon_-^2=-\partial_{\omega_1}\varepsilon_-=\sum_kk\,\omega_2^{-k-1}\omega_1^{k-1}$). With the fully
normal-ordered $\mathcal N=e^{A^+}e^{B^+}e^{A^-}e^{B^-}$ (multiplications commute, and so do differentiations):
\begin{equation}\tag{51}
e^{A^+}e^{A^-}e^{B^+}e^{B^-}=e^{\kappa}\mathcal N\ \text{in }\cL^+,\qquad e^{B^+}e^{B^-}e^{A^+}e^{A^-}=e^{\kappa'}\mathcal N\ \text{in }\cL^-.
\end{equation}
\emph{In which ring these identities hold.} Each side of (51), and of the two commutation rules before it, is a
formal power series in the $z_n,z'_n$ whose coefficients are $\mathrm{End}(V)$-valued formal series in
$\omega_1^{\pm1},\omega_2^{\pm1}$. Such a product is read as follows: fix a monomial in the $z,z'$, a power of
$\omega_2$ (resp.\ of $\omega_1$), and two monomials $t^\sigma,t^{\sigma'}$; then $[t^\sigma](\cdots)\,t^{\sigma'}$ is a
finite sum of Laurent monomials in the other variable, and an identity means that these finite sums agree for every
such choice. The finiteness is the argument of \cite[App.~B]{ETH}. In every product here $e^{A^+}$ stands left of
$e^{A^-}$ and $e^{B^+}$ left of $e^{B^-}$. If moreover $e^{B^-}$ does not stand left of $e^{A^+}$, the result lies in
$\cL^+$: (i) the differentiations from $B^-$ are applied before all multiplications, so their indices add up to at most
$|\sigma'|$; (ii) the $\omega_2$-exponent of a term $z'_n\nabla_{-k}$ of $B^+$ is $k-n\ge0$ and grows with $k$, while
those from $B^-$ are bounded by (i), so a fixed power of $\omega_2$ leaves finitely many choices for the
multiplications from $B^+$; (iii) with the $z$-monomial and all weights now bounded, the indices from $A^\pm$ are
bounded too. The $\omega_2$-exponents are bounded below, so the coefficient lies in $\cL^+$. Likewise the result lies
in $\cL^-$ (with $\omega_1,\omega_2$ exchanged) if $e^{A^-}$ does not stand left of $e^{B^+}$. In particular
$\kappa$, $e^\kappa$ and both sides of the first identity in (51) are in $\cL^+$, and $\kappa'$ and the second identity
in $\cL^-$.

\emph{One Laurent polynomial.} Fix $\sigma,\sigma'\in\cP_E$ and write $\ip X=[t^\sigma]X\,t^{\sigma'}$. Every coefficient
$n_{\mu'',\nu''}$ of $z^{\mu''}z'^{\nu''}$ in $\ip{\mathcal N}$ is in $\cL$ (only $\nabla_k$ with $|k|\le E$ contribute between
two weight-$E$ states), and $n_{\emptyset,\emptyset}=\delta_{\sigma\sigma'}$. The coefficient of $z^{\mu'}z'^{\nu'}$ in $e^\kappa$
is $\Gamma_{\mu',\nu'}\varepsilon_+^{|\mu'|+|\nu'|}$ for a number $\Gamma_{\mu',\nu'}$ (each $z_jz'_n$ carries $\varepsilon^{j+n}$), and in
$e^{\kappa'}$ it is $\Gamma_{\mu',\nu'}\varepsilon_-^{|\mu'|+|\nu'|}$. Comparing coefficients of $z^\mu z'^\nu$ in (51):
\[
f_\pm:=\ip{Y(a,\omega_1)Y(b,\omega_2)}\ \text{resp.}\ \ip{Y(b,\omega_2)Y(a,\omega_1)}
=\mu!\,\nu!\sum\Gamma_{\mu',\nu'}\,\varepsilon_\pm^{|\mu'|+|\nu'|}\,n_{\mu'',\nu''},
\]
summed over splittings $\mu=\mu'\sqcup\mu''$, $\nu=\nu'\sqcup\nu''$. Since $|\mu'|+|\nu'|\le H$, multiplying by
$(\omega_1-\omega_2)^H$ gives in both cases the same Laurent polynomial
$F=\mu!\nu!\sum\Gamma_{\mu',\nu'}(\omega_1-\omega_2)^{H-|\mu'|-|\nu'|}n_{\mu'',\nu''}\in\cL$, with
$F(1,1)=\mu!\nu!\,\Gamma_{\mu,\nu}\,\delta_{\sigma\sigma'}$ (every other term has a factor $\omega_1-\omega_2$).

\emph{Reading off.} Only weight-preserving products contribute to $\ip\cdot$, so
\[f_+=\sum_r\ip{a[r]b[-r]}\,\omega_1^{r-h'}\omega_2^{-r-h},\qquad f_-=\sum_r\ip{b[-r]a[r]}\,\omega_1^{r-h'}\omega_2^{-r-h},\]
with no terms for $r>E$ in $f_+$ ($b[-r]t^{\sigma'}\in V_{E-r}$) and none for $r<-E$ in $f_-$. Multiplying by
$(\omega_1-\omega_2)^H=\sum_i(-1)^i\binom Hi\omega_1^{H-i}\omega_2^i$ and collecting
$\omega_1^s\omega_2^{-s}$ ($r=s-h+i$):
\[F=\sum_sF_s\,\omega_1^s\omega_2^{-s},\qquad F_s=\ip{Q_s(a,b)}=\ip{Q'_s(a,b)}.\]
This is (a). For $s>E+h$, every $r=s-h+i\ge s-h>E$, so all terms of the first form vanish; for $s<-E-h$, every
$r=s-h+i\le s-h+H=s+h'<-E$, so all terms of the second form vanish ($a[r]t^{\sigma'}\in V_{E+r}=0$). This is (b).
Finally $\sum_sF_s=F(1,1)=\mu!\nu!\Gamma_{\mu,\nu}\delta_{\sigma\sigma'}$, and $\Gamma_{\mu,\nu}$ does not depend on $E$.
To identify it take $E=0$: then $\ip{\mathcal N}=1$ (the differentiations kill $1$, and the constant term of
$e^{A^+}e^{B^+}1$ is $1$), so $f_+=\mu!\nu!\Gamma_{\mu,\nu}\varepsilon_+^H$, whose coefficient of $\omega_2^0$ is
$\mu!\nu!\Gamma_{\mu,\nu}\,\omega_1^{-H}$. In $f_+=\sum_r\ip{a[r]b[-r]}\omega_1^{r-h'}\omega_2^{-r-h}$ the $\omega_2^0$ term is
$r=-h$, i.e.\ $a[-h]\,(b[h]1)\,\omega_1^{-H}$. And $b[h]1=b$: on $1$ only terms of (47) with all $k_i=-l_i<0$ survive,
$\sum l_i=h=\sum\nu_i$, and $\binom{l_i-1}{\nu_i-1}=0$ for $l_i<\nu_i$, so $l_i=\nu_i$ and the term is $t^\nu$. Hence
$\mu!\nu!\Gamma_{\mu,\nu}=a[-h]b$, which is (c).
\end{proof}

\begin{remark}\label{rem:locality}
For homogeneous $a$, (a) is the coefficient form of locality of the fields $Y(a,\omega_1),Y(b,\omega_2)$ of the
Heisenberg vertex algebra, $(\omega_1-\omega_2)^H[Y(a,\omega_1),Y(b,\omega_2)]=0$ \cite[\S2.1 and Prop.~2.4]{Fre},
\cite[Def.~1.5.1]{MN}, restricted to the weight-preserving products, as \cite[App.~H]{ETH} observes. A reader who
knows that theory can take (a) for homogeneous $a$ from it. It does not by itself give (b), (c), or the bookkeeping
for inhomogeneous $a$ through $a^{\langle i\rangle}$; those are proved above directly, and they are what the identity
\eqref{eq:cross} needs.
\end{remark}

\section{Counting}\label{sec:count}
Let $\kappa_{\rm HR}=\pi\sqrt{2/3}/\ln2\approx3.70$. By Hardy--Ramanujan, $\kappa_{\rm HR}\sqrt E-\log_2E-C\le\log_2p(E)\le\kappa_{\rm HR}\sqrt E+C$
(the form used in \cite[(46)]{ETH}). For $N=2^\lambda$ and $h$ least with $p(h)\ge N$: $p(h-1)<2^\lambda$ gives
$\kappa_{\rm HR}\sqrt h\le\lambda+O(\log\lambda)$ and $h=O(\lambda^2)$. Since $p$ is nondecreasing,
$K_h\le(2h+1)(4h+1)\,p(2h)$, and $\log_2p(2h)\le\sqrt2\,\kappa_{\rm HR}\sqrt h+C\le\sqrt2\,\lambda+O(\log\lambda)$. So
$K_h\le N^{\sqrt2}(\log N)^{O(1)}$. Exact values (computed by \texttt{check\_lemmas.py}):

\begin{center}
\begin{tabular}{rrrr}
\toprule
$\lambda$ & $h$ & $\log_2K_h$ & $\log_2K_h/\lambda$\\
\midrule
10 & 23 & 22.7 & 2.272\\
20 & 61 & 38.4 & 1.922\\
40 & 186 & 68.9 & 1.723\\
60 & 372 & 98.5 & 1.641\\
70 & 489 & 113.2 & 1.618\\
80 & 620 & 127.8 & 1.598\\
100 & 927 & 156.9 & 1.569\\
200 & 3361 & 300.9 & 1.505\\
400 & 12663 & 586.4 & 1.466\\
600 & 27849 & 870.9 & 1.451\\
\bottomrule
\end{tabular}
\end{center}
The exponent first drops below $1.6147$ at $\lambda=72$, and approaches $\sqrt2$ only like $\log\lambda/\lambda$. For small $N$
the identity uses far more than $N^2$ products ($K_4=307$ for $N\le5$, $K_5=707$ for $N\le7$): it is purely asymptotic.

\section{What is checked, and where}\label{sec:map}

{\footnotesize\raggedright
\noindent\begin{tabularx}{\textwidth}{@{}>{\raggedright\arraybackslash}p{3.3cm}>{\raggedright\arraybackslash}X@{}}
\toprule
Step & Lean declaration (paths relative to \texttt{eth/proof/rates/GeneralPieces/}; those beginning \texttt{statements/} or \texttt{rates/} are relative to \texttt{eth/proof/}) \\
\midrule
Definition 7, Theorem B (statements) & \texttt{Pieces.Family.Is} (\texttt{statements/PiecesStatement.lean}); \texttt{AppHalfPieces.ThmHalfPiecesPure}, \texttt{LemmaQ}, \texttt{LemmaDual}, \texttt{PieceFacts} (\texttt{statements/AppHalfPiecesStatement.lean})\\
(47): finiteness, linearity; (O1) & \texttt{GeneralPieces.opFacts} (\texttt{Ring/OpFacts.lean}); \texttt{GeneralPieces.o1} (\texttt{Ring/Op.lean})\\
(51), the two orders & \texttt{GeneralPieces.GenF.stTwoOrders} (\texttt{GenF/Commute.lean}); \texttt{GeneralPieces.GenF.stExpansion} (\texttt{GenF/Main.lean})\\
One Laurent polynomial $F$; $F(1,1)$; the constant & \texttt{Expansion.dW\_pow\_mul\_fP}, \texttt{Expansion.dW\_pow\_mul\_fM}, \texttt{Expansion.ev\_F} (\texttt{Series/OneLaurent.lean}); \texttt{Expansion.const} (\texttt{Series/Constant.lean}), all in \texttt{GeneralPieces.Series}\\
Lemma \ref{lem:Q} (a), (b), (c) & \texttt{GeneralPieces.Series.Qk\_eq\_Qk'\_of}, \texttt{Qk\_eq\_zero\_of}, \texttt{sum\_Qk\_of} (\texttt{Series/LemmaQ.lean}); assembled as \texttt{GeneralPieces.lemmaQ} (\texttt{Ring/LemmaQ.lean}) via \texttt{GeneralPieces.stBridge} (\texttt{Ring/Bridge.lean})\\
Lemma \ref{lem:dual}; $h!/z_\mu\in\Z$ & \texttt{GeneralPieces.Dual.lemmaDual}, with \texttt{theta\_eq}, \texttt{theta\_of\_lt}, \texttt{theta\_self}; \texttt{GeneralPieces.Dual.zDvd} (\texttt{Ring/Dual.lean})\\
\eqref{eq:cross}, (54), (55), integrality & \texttt{GeneralPieces.pieceFacts} (\texttt{Ring/Pieces.lean}); per piece \texttt{GeneralPieces.ThmHalfPieces.sum\_A}, \texttt{row\_A}, \texttt{col\_A} (\texttt{Ring/ThmHalfPieces.lean})\\
Counting, $h=O(\lambda^2)$ & \texttt{GeneralPieces.Bits.exists\_bits\_le}, \texttt{exists\_hOf\_le} (\texttt{Ring/Bits.lean}), which take (46) as a hypothesis; (46), with $B\log_2E$ in place of $\log_2E$, is \texttt{ClauseBound.Part.exists\_bounds} (\texttt{rates/ClauseBound/Part/Bounds.lean}), used in \texttt{thmHalfPiecesPure}\\
Theorem B & \texttt{GeneralPieces.thmHalfPiecesPure} (\texttt{Ring/ThmHalfPieces.lean}), applied to \texttt{AppHalfPieces.pieceFacts} (\texttt{Main.lean})\\
Corollary \ref{cor}, bound form over $\Q$, $N=2^\lambda$ & \texttt{hurwitzSqrt2} (our \texttt{HurwitzProof.lean})\\
\bottomrule
\end{tabularx}\par}
\medskip

\paragraph{Independent numerical checks.} \texttt{check\_lemmas.py} (Python, exact integer arithmetic, written
from the paper's text, not from \cite{Sie} or the Lean code) implements (47), (48), $a^\mu$, $b^\mu$ and checks:
the Heisenberg relations; Lemma~\ref{lem:dual} for all $\mu,\nu\in\cP_h$, $h\le5$; Lemma~\ref{lem:Q}(a), (b), (c) for
$a=a^\mu$ (inhomogeneous), $b=b^\nu$, $h\le4$, on every basis monomial of $V_E$, $E\le4$, for $|s|\le2(E+h)+2$ (for
$s>E+h$ the vanishing of $Q_s$ is immediate, but for $s<-(E+h)$, and for $Q'_s$ with $s>E+h$, it is a real cancellation;
the lemma for general $a,b$ is not tested numerically);
\eqref{eq:cross} on all $p(h)^4$ index tuples for $h\le5$; both factorizations term by term for $h\le4$; and the identity of
Theorem~\ref{thm:main}, multiplied out as a polynomial in $2p(h)$ variables, for $h=2,\dots,5$ ($K_h=40,120,307,707$).
All 44 checks pass (about 160 seconds on a laptop). \cite{Sie} reports the same identity checked for $N=4$ and $N=8$ with a different implementation; we reproduced its $N=4$ case here.

\paragraph{Credit and provenance.} The observation that Theorem~B of \cite{ETH} gives $N^{\sqrt2+o(1)}$ for
Hurwitz's problem was published first by J.~P.~Siemaszko on 7 November 2026 \cite{Sie}, including the translation
to \cite[Lemma~3.5]{HWY} and instance checks. The construction and the proofs above are those of \cite{ETH}, re-presented. This note, its
Lean bridge and its scripts were produced with Claude models (Anthropic) in sessions directed by the author, who
has not verified the mathematics independently of the checks described.

\begin{thebibliography}{9}
\bibitem{ETH} Claude (Anthropic), \emph{A $2^{o(n)}$ time algorithm for 3-SAT}, 26 Oct 2026, \url{https://github.com/anthropics/formal-math/tree/main/eth}.
\bibitem{FM} \texttt{anthropics/formal-math}, \texttt{eth/}, commit \texttt{16a5ce41304fcdaea4e0be7c132e1c361dc26da5}.
\bibitem{Fre} E.~Frenkel, \emph{Vertex algebras and algebraic curves}, S\'eminaire Bourbaki 1999/2000, expos\'e 875, Ast\'erisque 276 (2002), 299--339; arXiv:math/0007054 (numbering cited from this version, as in \cite{ETH}).
\bibitem{Hru} P.~Hrube\v{s}, \emph{A subquadratic upper bound on Hurwitz's problem and related non-commutative polynomials}, ECCC TR24-099, 2024.
\bibitem{HruCCC} P.~Hrube\v{s}, \emph{A subquadratic upper bound on sum-of-squares composition formulas}, CCC 2024, LIPIcs 300, 12:1--12:11; ECCC TR24-026 (theorem numbering as cited in \cite[App.~G]{ETH}).
\bibitem{HWY} P.~Hrube\v{s}, A.~Wigderson, A.~Yehudayoff, \emph{Non-commutative circuits and the sum-of-squares problem}, J. Amer. Math. Soc. 24 (2011); ECCC TR10-021.
\bibitem{MN} A.~Matsuo, K.~Nagatomo, \emph{Axioms for a vertex algebra and the locality of quantum fields}, MSJ Memoirs 4, 1999.
\bibitem{Sie} J.~P.~Siemaszko, \emph{An $N^{\sqrt2+o(1)}$ upper bound for Hurwitz's sum-of-squares problem from the ETH-refutation construction}, v0.5, Nov 2026, DOI 10.5281/zenodo.24024689.
\end{thebibliography}

\end{document}
